\documentclass{article}
\usepackage[utf8]{inputenc}
\usepackage[T1]{fontenc}
\usepackage{geometry}
\geometry{
    a4paper,
    left=2.5cm,
    right=2.5cm,
    top=2.5cm,
    bottom=2.5cm
}
\usepackage{graphicx}
\usepackage{float}
\usepackage{hyperref}
\usepackage[english]{babel}
\usepackage{amsmath}
\usepackage{amssymb}
\usepackage{array}
\usepackage{booktabs}
\usepackage{caption}
\usepackage{subcaption}
\usepackage{enumitem}
\usepackage{fancyhdr}
\usepackage{titlesec}
\usepackage{url}
\usepackage{textcomp}
\usepackage{xcolor}
\usepackage{listings}
\usepackage{multirow}
\usepackage{longtable}
\usepackage{tabularx}
\usepackage{times}

% Set up page style
\pagestyle{fancy}
\fancyhf{}
\rfoot{\thepage}
\renewcommand{\headrulewidth}{0pt}

% Title formatting
\titleformat{\chapter}[display]{\normalfont\bfseries}{\chaptertitle\thechapter}{0pt}{\Huge}
\titlespacing*{\chapter}{0pt}{50pt}{40pt}

% Document info
\title{Template Document}
\author{Academic Support Agent}
\date{\today}

\begin{document}
\maketitle

% %%%%%%%%%%%%%%%%%%%%%%%%%%%%%%%%%%
% %% Pagina de titulo
% %%%%%%%%%%%%%%%%%%%%%%%%%%%%%%%%%%

\begin{center}
\includegraphics[scale=0.2]{figuras/brasao.jpg} \\
PONTIF CIA UNIVERSIDADE CAT LICA DE MINAS GERAIS \\
Instituto de Ci ncias Exatas e de Inform tica

% \vspace{1.0cm}

\end{center}

 \vspace{0cm} {
 \singlespacing \Large{\monog \symbolfootnote[1]{Artigo apresentado ao Instituto de Ci ncias Exatas e Inform tica da Pontif cia Universidade Cat lica de \linebreak Minas Gerais, campus Contagem, como pr -requisito parcial para obten o do t tulo de Bacharel em Sistemas de Informa o.} \\ }
 % \normalsize{\monogES}
 }

\vspace{1.0cm}

\begin{flushright}
\singlespacing 
\normalsize{\AutorA \footnote{\funcaoA \cursA \,-- \emailA . }} \\
\normalsize{\AutorB \footnote{\funcaoB \cursB \,-- \emailB . }} \\
%\normalsize{\AutorC \footnote{\funcaoC \cursC \,-- \emailC . }} \\
%\normalsize{\AutorD \footnote{\funcaD \cursD \,-- \emailD . }} \\
\end{flushright}
\thispagestyle{empty}

\vspace{1.0cm}

\begin{abstract}
\noindent
O resumo   redigido em par grafo  nico. Um bom resumo deve conter uma frase para cada um dos t picos: contexto, problema, justificativa, objetivo, desenvolvimento, resultados e conclus o do documento. Conv m usar o verbo de forma impessoal. Seguindo o formato de artigo de peri dico, o resumo dever  conter de 100 a 250 palavras. Evite o uso de s mbolos e contra es que n o sejam de uso corrente, bem como o uso de f rmulas, equa es e refer ncias que n o sejam absolutamente necess rios. Quando seu emprego for imprescind vel, defini-los na primeira vez que aparecerem. Abaixo do resumo apresente a express o ``Palavras-chave'' e, ap s os dois pontos, os termos com as iniciais em letra min scula e separadas com ponto e v rgula (exceto para nomes pr prios e cient ficos) e finalizadas por ponto final. 
\\\textbf{\keyword{Palavras-chave:}} template; \LaTeX; Abakos; peri dicos.
\end{abstract}

%%%%%%%%%%%%%%%%%%%%%%%%%%%%%%%%%%%%%%%%%%%%%%%%%%%%%%%%%
\selectlanguage{english}
\begin{abstract}
\noindent
The abstract is written in a single paragraph. A good abstract should include the motivation, objective, methodology, results, and conclusion of the document. It is advisable to use impersonal verbs. Following the journal article format, the abstract should contain between 100 and 250 words. Avoid using symbols and contractions that are not commonly used, as well as formulas, equations, and references that are not absolutely necessary. When their use is essential, define them the first time they appear. Below the abstract, present the term ``Keywords'' and, after the colon, list the terms with lowercase initials separated by semicolons (except for proper and scientific names) and ending with a period.
\\\textbf{\keyword{Keywords:}} template; \LaTeX; Abakos; periodics.
\end{abstract}

\selectlanguage{brazilian}
 \onehalfspace  % espa amento 1.5 entre linhas
 \setlength{\parindent}{1.25cm}

%%%%%%%%%%%%%%%%%%%%%%%%%%%%%%%%%%%%%%%%%%%%%%%%%
%% INICIO DO TEXTO
%%%%%%%%%%%%%%%%%%%%%%%%%%%%%%%%%%%%%%%%%%%%%%%%%

% Content from textos would be included here


%%%%%%%%%%%%%%%%%%%%%%%%%%%%%%%%%%%
%% FIM DO TEXTO
%%%%%%%%%%%%%%%%%%%%%%%%%%%%%%%%%%%

% \selectlanguage{brazil}
%%%%%%%%%%%%%%%%%%%%%%%%%%%%%%%%%%%
%% Inicio bibliografia
%%%%%%%%%%%%%%%%%%%%%%%%%%%%%%%%%%%

 \newpage
\singlespace{
\renewcommand\refname{REFER NCIAS}

}

%Inclus o do arquivo abntex2-alf.bst como solu o para adequa o   ABNT NBR 10520:2023 quanto  s cita es, que n o s o mais em caixa-alta
\end{document}